\documentclass[../DoAn.tex]{subfiles}
\begin{document}

Chương này trình bày bốn đóng góp kỹ thuật nổi bật nhất của đề tài: (1) phương án tích hợp GPS và kết nối di động trên một module phần cứng duy nhất, (2) thuật toán geofencing hai chế độ hỗ trợ cả vùng an toàn hình tròn lẫn vùng an toàn dạng hành lang đường đi, (3) kiến trúc truyền thông lai kết hợp MQTT và WebSocket, và (4) cơ chế cảnh báo đa kênh song song. Mỗi đóng góp được trình bày theo ba phần: bài toán đặt ra, giải pháp thiết kế, và kết quả đạt được.

% ===========================================================
\section{Tích hợp GPS và kết nối di động trên một module phần cứng}
\label{section:5.1}

\subsection{Bài toán và thách thức}

Yêu cầu cốt lõi của hệ thống là thiết bị phải hoạt động \textbf{hoàn toàn độc lập}, không phụ thuộc điện thoại hay hạ tầng Wi-Fi, trong khi vẫn đảm bảo cập nhật vị trí liên tục theo thời gian thực. Điều này đồng nghĩa với việc một thiết bị nhỏ gọn phải tích hợp đồng thời hai chức năng độc lập: \textit{định vị vệ tinh} (GPS) và \textit{truyền dữ liệu qua mạng di động} (GPRS/NB-IoT).

Cách tiếp cận truyền thống là sử dụng hai module riêng biệt: một module GPS (ví dụ u-blox NEO-6M) và một modem di động (ví dụ SIM800L), mỗi module giao tiếp qua một cổng UART riêng. Phương án này làm tăng đáng kể số lượng linh kiện trên bo mạch, cần hai anten riêng biệt (GPS patch antenna và GSM wire antenna), và đòi hỏi quản lý nguồn độc lập cho từng module. Quan trọng hơn, firmware phải điều phối đồng thời hai luồng giao tiếp UART hoàn toàn khác nhau với logic khởi động và khởi động lại riêng rẽ — làm tăng đáng kể độ phức tạp mã nguồn và nguy cơ xung đột tài nguyên.

\subsection{Giải pháp thiết kế}

Đề tài lựa chọn module SIM7600 (được giới thiệu chi tiết tại mục \ref{subsection:3.1.2}) — module tích hợp cả bộ thu GNSS (GPS/GLONASS/BeiDou) lẫn modem di động 4G LTE Cat-1 trong một linh kiện duy nhất với một anten duy nhất và một cổng UART duy nhất.

Đóng góp không nằm đơn thuần ở việc chọn đúng linh kiện, mà ở \textbf{thiết kế tích hợp phần mềm điều khiển}. Thư viện TinyGSM trừu tượng hóa toàn bộ tập lệnh AT command của SIM7600 thành API cấp cao, cho phép ESP32 gọi \texttt{modem.getGPS()} để lấy tọa độ và \texttt{client.publish()} để gửi MQTT — tất cả thông qua cùng một kênh UART2, không cần chuyển đổi bus vật lý.

Thách thức phát sinh là GPS cần thời gian khởi động lạnh khoảng 30 giây để bắt vệ tinh, trong khi đó kết nối GPRS phải được duy trì liên tục. Firmware được thiết kế theo \textbf{máy trạng thái} (state machine) với các trạng thái: \texttt{INIT} → \texttt{GSM\_CONNECT} → \texttt{MQTT\_CONNECT} → \texttt{GPS\_WAIT\_FIX} → \texttt{TRACKING}. Modem khởi động kết nối di động trước, sau đó song song bật GPS; tọa độ chỉ được gửi khi cả hai đều sẵn sàng. Nếu kết nối mạng bị mất, firmware chuyển về trạng thái \texttt{GSM\_CONNECT} để thử kết nối lại mà không làm gián đoạn phiên GPS đang hoạt động.

\subsection{Kết quả đạt được}

So với phương án hai module riêng biệt, giải pháp trên đạt được:
\begin{itemize}
    \item Giảm số linh kiện chính trên bo mạch, đơn giản hóa sơ đồ nguyên lý và bố cục PCB.
    \item Một anten duy nhất phục vụ cả GPS lẫn di động, giảm chi phí và kích thước thiết bị.
    \item Firmware chỉ cần quản lý một giao diện UART, giảm nguy cơ tranh chấp tài nguyên và đơn giản hóa quá trình gỡ lỗi.
    \item Máy trạng thái đảm bảo thiết bị tự phục hồi sau khi mất tín hiệu GPS hoặc rớt mạng mà không cần can thiệp thủ công.
\end{itemize}

% ===========================================================
\section{Thuật toán geofencing hai chế độ}
\label{section:5.2}

\subsection{Bài toán và thách thức}

Geofencing truyền thống xác định vùng an toàn bằng một hình tròn với tâm và bán kính cố định. Cách tiếp cận này phù hợp khi trẻ ở trong một khu vực cụ thể như sân trường hay công viên, nhưng thiếu linh hoạt với thực tế sử dụng hàng ngày: trẻ thường di chuyển theo \textbf{tuyến đường cố định} từ nhà đến trường, đến trung tâm học thêm, hoặc đến nhà ông bà. Vùng an toàn dạng hành lang đường đi (path-based corridor) tự nhiên và chính xác hơn nhiều so với hình tròn bán kính lớn bao trùm cả khu vực rộng không cần thiết.

Ngoài ra, bài toán định nghĩa ``vị trí trẻ có nằm trong vùng an toàn không'' đòi hỏi tính toán khoảng cách địa lý chính xác trên bề mặt Trái Đất — không phải khoảng cách Euclid trên mặt phẳng, vì tọa độ GPS là kinh độ và vĩ độ trên hình cầu.

\subsection{Giải pháp thiết kế}

Mỗi đối tượng geofence trong collection \texttt{geofence\_configs} lưu một tâm mặc định (\texttt{centerLat}, \texttt{centerLng}), bán kính \texttt{radiusM}, và một mảng \texttt{path} tùy chọn chứa danh sách điểm tọa độ của tuyến đường. \textbf{Hình dạng vùng an toàn được xác định hoàn toàn bởi nội dung của trường \texttt{path}}:

\textbf{a. Vùng hình tròn (\texttt{path} rỗng):} Backend tính khoảng cách từ tọa độ hiện tại $(lat_p, lng_p)$ đến tâm geofence $(lat_c, lng_c)$ bằng công thức Haversine (trình bày tại mục \ref{section:3.5}). Nếu khoảng cách $d > radiusM$, hệ thống kích hoạt cảnh báo. Đây là lựa chọn phù hợp khi phụ huynh muốn giới hạn trẻ trong một khu vực cụ thể như sân trường hay công viên.

\textbf{b. Vùng hành lang đường đi (\texttt{path} có $\geq 2$ điểm):} Phụ huynh vẽ tuyến đường trên bản đồ trong ứng dụng; tuyến đường được lưu dưới dạng danh sách các điểm tọa độ $\{(lat_1, lng_1), (lat_2, lng_2), \ldots, (lat_n, lng_n)\}$ trong trường \texttt{path} của geofence. Backend tính khoảng cách từ điểm vị trí hiện tại đến \textbf{từng đoạn thẳng} của tuyến đường theo phép chiếu vuông góc:

Với mỗi đoạn thẳng $\overrightarrow{AB}$, điểm chiếu $P'$ của điểm $P$ lên đoạn thẳng được tìm bằng tham số $t$:
\begin{equation}
t = \frac{\overrightarrow{AP} \cdot \overrightarrow{AB}}{|\overrightarrow{AB}|^2}, \quad t \in [0, 1]
\label{eq:projection}
\end{equation}

Nếu $t < 0$, điểm gần nhất trên đoạn là $A$; nếu $t > 1$, điểm gần nhất là $B$; ngược lại điểm gần nhất là $P' = A + t \cdot \overrightarrow{AB}$. Khoảng cách Haversine từ $P$ đến điểm gần nhất này được tính và so sánh với bán kính hành lang \texttt{radiusM}. Khoảng cách nhỏ nhất trong tất cả các đoạn là $d_{min}$; nếu $d_{min} > radiusM$, cảnh báo được kích hoạt.

Độ phức tạp thuật toán là $O(n \cdot m)$ với $n$ là số geofence đang hoạt động và $m$ là số đoạn thẳng trong tuyến đường — đủ hiệu quả cho quy mô nguyên mẫu và thời gian xử lý không đáng kể so với latency mạng.

Cả hai hình dạng được xử lý trong \textbf{cùng một pipeline backend}: mỗi khi nhận được tọa độ mới từ MQTT, hàm kiểm tra geofence chạy lần lượt qua tất cả geofence đang hoạt động của thiết bị, kiểm tra độ dài \texttt{path} để chọn nhánh tính khoảng cách tương ứng (tròn hoặc hành lang), và lưu kết quả \texttt{insideGeofence} cùng \texttt{distanceFromCenterM} vào bản ghi \texttt{locations}. Việc suy ra hình dạng trực tiếp từ dữ liệu \texttt{path} — thay vì từ một cờ chế độ riêng — giúp loại bỏ khả năng hai khái niệm rơi vào trạng thái không nhất quán với nhau.

Song song với lựa chọn hình dạng, hệ thống còn có trường \texttt{mode} (\texttt{fixed}/\texttt{mobile}) trên mỗi geofence, cho phép tâm của vùng hoặc cố định do phụ huynh đặt (\texttt{fixed}), hoặc tự động bám theo vị trí GPS hiện tại của điện thoại phụ huynh (\texttt{mobile}, đồng bộ định kỳ mỗi 10 giây). Đây là một trục lựa chọn độc lập với hình dạng: một vùng GPS-bám-theo luôn ở dạng hình tròn (bán kính cố định quanh vị trí tức thời), trong khi vùng hành lang đường đi chỉ áp dụng cho geofence có tâm cố định.

\subsection{Kết quả đạt được}

Giải pháp geofencing với hai hình dạng (tròn và hành lang) mang lại tính linh hoạt đáng kể trong thực tế sử dụng. Phụ huynh có thể kết hợp cả hai loại trong cùng một hệ thống: ví dụ đặt vùng hình tròn xung quanh trường học và vùng hành lang dọc theo con đường trẻ thường đi. Thuật toán hoàn toàn tự chủ trên backend, không phụ thuộc dịch vụ bên ngoài, cho phép kiểm soát hoàn toàn logic phát hiện vi phạm và dễ dàng mở rộng thêm hình dạng mới trong tương lai.

% ===========================================================
\section{Kiến trúc truyền thông lai MQTT–WebSocket}
\label{section:5.3}

\subsection{Bài toán và thách thức}

Hệ thống phải truyền dữ liệu vị trí từ thiết bị nhúng đến ứng dụng di động của phụ huynh theo thời gian thực, qua hai đường truyền có đặc điểm hoàn toàn khác nhau: \textbf{đường truyền IoT} (thiết bị nhúng – mạng di động GPRS – server) và \textbf{đường truyền di động} (server – Internet – smartphone).

Không có giao thức đơn lẻ nào tối ưu cho cả hai đường truyền:
\begin{itemize}
    \item MQTT rất phù hợp cho đường truyền IoT: nhẹ, hỗ trợ QoS, pub/sub tách biệt publisher và subscriber, chịu kết nối không ổn định tốt. Tuy nhiên MQTT không phải chuẩn native của nền tảng di động; kết nối MQTT trực tiếp từ Flutter đòi hỏi quản lý phiên phức tạp và không tận dụng được cơ chế push notification của hệ điều hành.
    \item WebSocket là chuẩn được hỗ trợ đầy đủ trên nền tảng Android/iOS, phù hợp cho kênh server-to-client realtime. Tuy nhiên WebSocket không có cơ chế QoS hay retry tích hợp, và không phù hợp để thiết bị nhúng có tài nguyên hạn chế sử dụng làm kênh chính.
    \item HTTP polling là phương án đơn giản nhất nhưng tạo độ trễ phụ thuộc vào chu kỳ polling và tốn tài nguyên cả phía client lẫn server khi tần suất polling cao.
\end{itemize}

\subsection{Giải pháp thiết kế}

Đề tài thiết kế kiến trúc \textbf{lai MQTT–WebSocket} với backend đóng vai trò cầu nối (bridge) giữa hai thế giới:

\begin{enumerate}
    \item \textbf{Tầng thiết bị → Backend:} ESP32 publish dữ liệu GPS lên MQTT broker (EMQX public broker, mục \ref{subsection:3.2.2}) theo topic phân cấp. Backend (FastAPI) đăng ký nhận toàn bộ topic qua Paho-MQTT subscriber chạy trong luồng nền riêng.

    \item \textbf{Xử lý nội bộ Backend:} Khi Paho-MQTT callback nhận được bản tin mới (chạy trong Python thread riêng), dữ liệu được đẩy vào \texttt{asyncio.Queue} để chuyển giao an toàn sang event loop bất đồng bộ của FastAPI. Đây là điểm kỹ thuật quan trọng: Paho-MQTT callback là đồng bộ, trong khi WebSocket broadcast là bất đồng bộ; hàng đợi \texttt{asyncio.Queue} là cơ chế an toàn để bridge giữa thread và coroutine.

    \item \textbf{Tầng Backend → Ứng dụng di động:} Một coroutine nền liên tục đọc từ hàng đợi và broadcast vị trí mới đến tất cả kết nối WebSocket đang hoạt động của thiết bị tương ứng. Ứng dụng Flutter duy trì một kết nối WebSocket duy nhất và cập nhật bản đồ ngay khi nhận được dữ liệu.
\end{enumerate}

Thiết kế này đạt được \textbf{phân tách trách nhiệm rõ ràng} giữa các tầng: tầng thiết bị chỉ cần biết địa chỉ MQTT broker; tầng ứng dụng chỉ cần biết địa chỉ WebSocket của backend; backend là thành phần duy nhất hiểu cả hai giao thức. Khi cần thay đổi broker MQTT hay mở rộng số client WebSocket, chỉ cần sửa cấu hình backend mà không ảnh hưởng đến firmware hay ứng dụng di động.

\subsection{Kết quả đạt được}

Kiến trúc lai đạt độ trễ end-to-end (từ khi ESP32 đọc tọa độ GPS đến khi bản đồ Flutter cập nhật) dưới 2 giây trong điều kiện mạng bình thường, trong đó phần lớn độ trễ đến từ chu kỳ cập nhật GPS 5 giây của phần cứng. Độ trễ truyền thông thuần túy (MQTT broker → FastAPI → WebSocket → Flutter) đo được dưới 300 mili giây trong điều kiện thử nghiệm. Kiến trúc này cũng cho phép nhiều phụ huynh (nhiều kết nối WebSocket) cùng theo dõi một thiết bị mà không cần thay đổi firmware hay logic MQTT.

% ===========================================================
\section{Cơ chế cảnh báo đa kênh song song}
\label{section:5.4}

\subsection{Bài toán và thách thức}

Cảnh báo geofence và tín hiệu SOS là hai sự kiện quan trọng nhất của hệ thống; phụ huynh phải nhận được ngay lập tức, kể cả khi không đang mở ứng dụng. Nếu chỉ dựa vào WebSocket, cảnh báo sẽ mất hoàn toàn khi ứng dụng đang chạy nền hoặc màn hình điện thoại tắt. Nếu chỉ dựa vào push notification của hệ điều hành (FCM/APNs), cảnh báo có thể bị trễ hoặc thất lạc do cơ chế quản lý pin của Android và iOS. Không kênh nào đơn lẻ đảm bảo được độ tin cậy yêu cầu cho tình huống khẩn cấp.

\subsection{Giải pháp thiết kế}

Khi backend phát hiện vi phạm geofence hoặc nhận tín hiệu SOS, hệ thống kích hoạt \textbf{hai kênh thông báo song song}, không phụ thuộc nhau:

\begin{enumerate}
    \item \textbf{WebSocket (trong ứng dụng):} Bản tin cảnh báo được broadcast đến tất cả kết nối WebSocket của phụ huynh đang xem bản đồ theo thời gian thực. Ứng dụng Flutter hiển thị thông báo popup và âm thanh cảnh báo ngay lập tức. Kênh này có độ trễ thấp nhất nhưng chỉ hoạt động khi ứng dụng đang mở.

    \item \textbf{Telegram Bot:} Backend gọi Telegram Bot API để gửi tin nhắn văn bản kèm tọa độ GPS và liên kết Google Maps trực tiếp đến số điện thoại/tài khoản Telegram của phụ huynh. Kênh này hoạt động hoàn toàn độc lập với ứng dụng di động; Telegram có cơ chế push notification riêng của mình và hoạt động tốt kể cả khi ứng dụng đang đóng. Với tín hiệu SOS, tin nhắn Telegram bao gồm thêm link bản đồ Google Maps định vị tức thì vị trí của trẻ.
\end{enumerate}

Đồng thời, bất kể kênh nào được kích hoạt, sự kiện cảnh báo đều được ghi vào collection \texttt{alerts} trong MongoDB với đầy đủ thông tin: loại cảnh báo, tọa độ, thời gian, tên geofence vi phạm (nếu có). Điều này cho phép phụ huynh xem lại toàn bộ lịch sử cảnh báo trong ứng dụng sau đó.

\subsection{Kết quả đạt được}

Cơ chế hai kênh song song đảm bảo phụ huynh nhận được cảnh báo trong mọi tình huống sử dụng thực tế: đang xem ứng dụng, hoặc đang dùng điện thoại cho việc khác nhưng vẫn nhận được qua Telegram. Trong thử nghiệm thực tế, cả hai kênh đều gửi thành công trong vòng 3 giây sau khi phát hiện vi phạm, với Telegram là kênh nhanh nhất (thường dưới 1 giây).

\end{document}
